\documentclass[hidelinks, headlogo, framedcover]{tfeEPCC}
\usepackage[spanish, es-tabla]{babel} % Comentar si se va a realizar en inglés
%\usepackage{parskip} % Quitar comentario si se quieren distinguir los párrafos con saltos y quitar las sangrías

\usepackage{pgfgantt}
\usepackage{tabularx}
\usepackage{xcolor}
%\usepackage{lipsum} % Eliminado para la versión final de impresión

\usepackage{tikz}
\usetikzlibrary{positioning, fit, backgrounds, shapes, arrows.meta, calc, babel, automata, er, matrix}
\usepackage{pgf-umlsd}
\usepackage{graphicx}
\usepackage{dirtree}
\usepackage{adjustbox}
\usepackage{multirow}
\usepackage{pgfplots}
\pgfplotsset{compat=1.18}

\newcommand{\mycomment}[1]{\hfill \textcolor{gray}{\textit{\small #1}}}

\definecolor{ontime}{RGB}{66, 133, 244}    % Azul - a tiempo
\definecolor{early}{RGB}{52, 168, 83}      % Verde - anticipado
\definecolor{late}{RGB}{234, 67, 53}       % Rojo - retrasado
\definecolor{inprogress}{RGB}{251, 188, 4} % Amarillo - en progreso

% Rellenar con los datos del estudiante y el trabajo
\estudiante{Alejandro Barrena Millán}{Grado}{Grado en Ingeniería Informática en Ingeniería del Software}
\title{Diseño y análisis de un agente LLM autónomo para auditorías de ciberseguridad en entornos aislados}
\keywords{Agente autónomo\and LLM\and Ciberseguridad\and Pentesting\and Human-in-the-Loop\and Docker\and Auditoría automatizada}

\begin{document}
\sloppy 
\pagebreak 

%%%%%%%%%%%%% PORTADA %%%%%%%%%%%%%
% Poner convocatoria y año de presentación
\portadaTFE{Febrero, 2025}

\pagestyle{empty}


%%%%%%%%%% CONTRAPORTADA %%%%%%%%%%
% Campos, separar con comas si se ponen varios: {
% studentGender
% tutor (obligatorio)
% tutorGender
% cotutor
% cotutorGender }
% Especificar un co-tutor lo añade a la contraportada
% Poner la letra f en cualquier campo de género añade una 'a' al título correspondiente (por ejemplo, studentGender = f hace que se muestre Autora)
\contraportadaTFE{tutor = Pablo García Rodríguez, cotutor = Alberto Fernández Isabel (URJC)}

%%%% ALEJANDRO BARRENA COSAS PERSONALIZADAS %%%%
\setlength{\parskip}{10pt} % Añade un pequeño espacio extra entre párrafos


%%%% CAPÍTULOS INTRODUCTORIOS %%%%
\pagenumbering{gobble}
%\chapter*{Agradecimientos} % Apartado opcional, descomentar si se quiere añadir
% Dedicatorias del trabajo, pueden ser en español o inglés dependiendo de lo que se considere más apropiado
%%%%%%%%%%%%%%%%%%%%%%%%%%%%%%%%%%%%%%%%%%%%%%%%%%%%
% DEDICATORIA
%%%%%%%%%%%%%%%%%%%%%%%%%%%%%%%%%%%%%%%%%%%%%%%%%%%%

%%%%%%%%%%%%%%%%%%%%%%%%%%%%%%%%%%%%%%%%%%%%%%%%%%%%
% AGRADECIMIENTOS
%%%%%%%%%%%%%%%%%%%%%%%%%%%%%%%%%%%%%%%%%%%%%%%%%%%%
\cleardoublepage
\newpage
\thispagestyle{empty}
\pagenumbering{gobble} 

% Truco 1: Reducir el espacio muerto antes del título
\begin{center}
    \vspace*{-2cm} % Sube el título hacia arriba
    \Huge \textbf{Agradecimientos}
    \vspace*{1cm}  % Espacio tras el título
\end{center}

% Truco 2: Ganar un par de líneas extra al final de la página
\enlargethispage{2cm}

\noindent En primer lugar, quiero expresar mi más profundo agradecimiento a mis padres, por todas las oportunidades que me han brindado y por haberme formado como la persona que soy hoy, apoyándome incondicionalmente en cada decisión. A mi madre, por ser quien más me ha apoyado día tras día, dándome la fuerza necesaria para no rendirme nunca. A mi padre, por haberme transmitido siempre su cariño e ilusión; hoy más que nunca podrá decir, por fin, que soy "su ingeniero".

\vspace{0.4cm} % Espacio controlado entre bloques

\noindent A toda mi familia, por estar siempre ahí. Especialmente a mi abuela Pili y a mi tía Mari, por su cariño inmenso y por haber sido un apoyo constante en cada etapa de mi vida. También quiero agradecer de forma muy especial a mi tío Manolo. Gracias por ayudarme a remontar el vuelo tras aquel difícil primer año; gracias por enseñarme a volver a aprender y por demostrarme que un mal comienzo no define el final del camino si se tiene la constancia necesaria para levantarse.

\vspace{0.4cm}

\noindent Por último, quiero dar las gracias a todos los amigos y compañeros con los que he compartido estos años de estudio. Hemos formado un auténtico equipo, afrontando cada examen como una batalla conjunta y apoyándonos mutuamente para hacer este trayecto mucho más ameno. En especial a Lete y Natera, por vuestro apoyo constante durante todo el desarrollo y por esa paciencia infinita al no quejaros cuando Robe traspasaba las paredes de mi habitación durante el desarrollo del trabajo; gracias por ser un apoyo fundamental en cada momento de bloqueo. También a Elena, por estar siempre ahí, escucharme en todo momento y por su apoyo incondicional a lo largo de todo el proceso. Gracias a todos por estar ahí, de verdad.

\chapter*{Dedicatoria}
\thispagestyle{empty}
\pagenumbering{gobble}

\begin{flushright}
    \textit{A mi abuelo Antonio, con quien tanto compartí \\ 
    y de quien guardo los mejores recuerdos. \\ 
    Gracias por contagiarme esa curiosidad \\ 
    y ese ingenio que llevamos los Barrena; \\ 
    casi sin darte cuenta, me dejaste esa forma de ver las cosas \\ 
    que hoy me ha traído hasta aquí.}

    \vspace{0.8cm}

    \textit{A mis abuelos Francisco y Santa, por haberme cuidado tanto \\ 
    y por ser una parte tan importante de mi infancia. \\ 
    Abuela, gracias por tu cariño inmenso que siempre llevo conmigo. \\ 
    Abuelo, todavía resuena tu voz en mi cabeza, \\ 
    con aquel acento granaino, diciendo ese “estuuudia” \\ 
    que tanto me repetías; espero que estés orgulloso de este resultado.}

    \vspace{0.8cm}

    \textit{Gracias por todo lo que habéis dejado en mí; \\ 
    vuestra huella es mi mejor guía.}
\end{flushright}

\cleardoublepage
\chapter*{Resumen}
Este Trabajo de Fin de Grado presenta el diseño, implementación y evaluación de un sistema agéntico autónomo basado en LLMs orientado a la realización de auditorías de ciberseguridad. La plataforma, se articula mediante una arquitectura de microservicios compuesta por tres componentes: un sandbox aislado de Linux para la ejecución segura de herramientas de pentesting, un backend orquestador que gestiona la comunicación con múltiples proveedores de inteligencia artificial, y un frontend web que ofrece una interfaz de supervisión humana con sistema de aprobación de comandos peligrosos.

El sistema ha sido validado a través de auditorías contra tres entornos de vulnerabilidad controlados, evaluando 8 modelos de lenguaje en tres niveles de infraestructura: hardware doméstico, computación de alto rendimiento (COMPUTAEX) y APIs comerciales en la nube. Los resultados demuestran que los modelos cloud de última generación alcanzan una gran eficacia en escenarios de pentesting, mientras que los modelos locales presentan limitaciones significativas asociadas a la capacidad de contexto y razonamiento. El mecanismo Human-in-the-Loop implementado clasifica comandos por nivel de riesgo con una precisión del 92.8\%, permitiendo una supervisión focalizada sin comprometer la fluidez operativa del agente.

El código fuente del proyecto se encuentra disponible en \url{https://github.com/Alexeido/TryHackYou}.

\noindent\textbf{Palabras clave}: \printkeywords.

\cleardoublepage
\chapter*{Abstract}
This Bachelor's Thesis presents the design, implementation and evaluation of an autonomous agentic system based on Large Language Models (LLMs) aimed at performing cybersecurity audits. The platform, named TryHackYou, is built on a microservices architecture comprising three components: an isolated sandbox based on Kali Linux for the secure execution of pentesting tools, a backend orchestrator built with FastAPI that manages communication with multiple AI providers, and a frontend developed in Next.js that provides a human supervision interface with a dangerous command approval system.

The system has been validated through 14 audits against three controlled vulnerability environments (DVWA, Metasploitable2 and OWASP Juice Shop), evaluating 8 language models across three infrastructure levels: consumer hardware, high-performance computing (COMPUTAEX) and commercial cloud APIs. Results demonstrate that state-of-the-art cloud models achieve 100\% effectiveness in solving pentesting scenarios, while local models exhibit significant limitations related to context capacity and reasoning ability. The implemented Human-in-the-Loop mechanism classifies commands by risk level with 92.8\% accuracy, enabling focused supervision without compromising the agent's operational fluency.

The project source code is available at \url{https://github.com/Alexeido/TryHackYou}.

\noindent\textbf{Keywords}: Autonomous agent, LLM, Cybersecurity, Pentesting, Human-in-the-Loop, Docker, Automated auditing.

% Índices con enlaces de las secciones, tablas y figuras
\cleardoublepage
\newpage
\pagestyle{plain}
\pagenumbering{Roman}
\setcounter{page}{1}
\setcounter{tocdepth}{1}
\tableofcontents
%\pagebreak
\cleardoublepage
\listoftables
\cleardoublepage
\listoffigures
\cleardoublepage


%%%%%%% ESTRUCTURA DEL TFG %%%%%%%%%%%%%%
% A.PORTADA (según la estructura indicada a continuación) 
% B.CONTRAPORTADA (según la estructura indicada a continuación) 
% C.ÍNDICE GENERAL DE CONTENIDOS  
% D.ÍNDICE DE TABLAS  
% E.ÍNDICE DE FIGURAS 
% F.RESUMEN (Podrá incluirse también en inglés, si así lo indica el Tutor1) 
% G.CUERPO DEL TRABAJO (según la estr
% uctura indicada a continuación) 
% H.REFERENCIAS BIBLIOGRÁFICAS 
% (Según norma ISO690)
% I.ANEXOS, si los hubiera
%%%%%%%%%%%%%%%%%%%%%%%%%%%%%%%%%%%%%%%
%%%%%% Cuerpo del trabajo
% 1.INTRODUCCIÓN 
% 2.OBJETIVOS 
% 3.ANTECEDENTES / ESTADO DEL ARTE 
% 4.MÉTODOLOGÍA 
% 5.IMPLEMENTACIÓN Y DESARROLLO (Cuando proceda) 
% 6.RESULTADOS Y DISCUSIÓN 
% 7.CONCLUSIONES 
%%%%%%%%%%%%%%%%%%%%%%%%%%%%%%%%%%%%%%%%%%%
%%%%%%%%%%%%%%%%%%%%%%%%%%%%%%%%%%%%%%%%%%%
% Tamaño: Normalizado UNE A-4, salvo planos. 
% Tipo y tamaño de letra del texto: Times New Roman 12 pt, Arial 12 pt o similar. 
% Interlineado del texto: 1,5 líneas. 
% Márgenes del texto: Superior, Inferior y Derecha, 2,5 cm; Izquierda, 4 cm. 
% Numeración de páginas en margen inferior derecha y tamaño 8 pt. 
% Las  figuras  serán  numeradas  y  tituladas  debajo  de  las  mismas  (indicando  su  fuente  si  no  son  de  elaboración  
% propia). 
% Las  tablas  serán  numeradas  y  tituladas  encima  de  las  mismas  (indicando  su  fuente  si  no  son  de  elaboración  
% propia). 
%%%%%%%%%%%%%%%%%%%%%%%%%%%%%%%%%%%%%%%%%%%
%%%%%%%%%%%%%%%%%%%%%%%%%%%%%%%%%%%%%%%%%%%
%%%%%%%%%%%%%%%%%%%%%%%%%%%%%%%%%%%%%%%%%%%
%%%%%%%%%%%%%%%%%%%%%%%%%%%%%%%%%%%%%%%%%%%


\cleardoublepage
\pagenumbering{arabic}
\pagestyle{fancy}
\setcounter{page}{1}

%%%%%%%% INTRODUCCIÓN %%%%%%%%
\import{chapters}{01_introduccion.tex}

%%%%%%%% OBJETIVOS %%%%%%%%
\import{chapters}{02_objetivos.tex}

%%%%%%%% ESTADO DEL ARTE %%%%%%%%
\import{chapters}{03_estado_del_arte.tex}

%%%%%%%% METODOLOGÍA %%%%%%%%
\import{chapters}{04_metodologia.tex}

%%%%%%%% IMPLEMENTACIÓN Y DESARROLLO %%%%%%%%
\import{chapters}{05_implementacion_y_desarrollo.tex}

%%%%%%%% RESULTADOS %%%%%%%%
\import{chapters}{06_resultados.tex}

%%%%%%%% CONCLUSIONES Y TRABAJOS FUTUROS %%%%%%%%
\import{chapters}{07_conclusiones_y_trabajos_futuros.tex}


% Si no se van a incorporar anexos, comentar el comando '\startappendices' junto a todos los imports de la carpeta 'appendices'
%%%%%%%%%%%%%%%%%%%%%%%%%%%%%%%%%%
%%%%%%%%%%% ANEXOS %%%%%%%%%%%%%%%
%%%%%%%%%%%%%%%%%%%%%%%%%%%%%%%%%%
%\startappendices

%%%%%%%% ANEXO 1 %%%%%%%%
% \import{appendices}{01_ejemplo_de_anexo.tex}

\cleardoublepage
\pagebreak
%%%%%%%%%%%%%%%%%%%%%%%%%%%%%%%%%%
%%%%%%%%%% AL FINAL %%%%%%%%%%%%%%
%%%%%%%%%%%%%%%%%%%%%%%%%%%%%%%%%%
\pagestyle{empty}
%%%% https://en.wikibooks.org/wiki/LaTeX/Bibliography_Management
%%%% https://www.bibtex.com
%%%%%%%% BIBLIOGRAFÍA %%%%%%%%
% Se puede cambiar el estilo de la bibliografía, por defecto es 'unsrturl'
% Por ejemplo: \makebibliography{alphaurl}
\makebibliography
\end{document}